% assets/a3-nullspace-slide.svg — null space 보정은 두 힘을 반대로 미끄러뜨린다
%
%   bash _tools/tikz2svg.sh assets/src/a3-nullspace-slide.tex assets/a3-nullspace-slide.svg
%
% 이 그림이 말하는 것
%   (가) 보정 전 — 우현 추진기의 방향각이 +-45 도 쐐기를 넘는다
%   (나) 보정 후 — 두 추진기의 F_y 를 +lambda / -lambda 로 밀면 쐐기 안으로 들어온다
%        요구 렌치는 하나도 바뀌지 않는다
%   (다) lambda 를 1 차 부등식 둘로 풀어 고르는 과정
%
% 화면 좌표 — 평면도 (위에서 내려다본 그림)
%   화면 x = y_b (우현 쪽),  화면 y = x_b (선수 쪽)
%   따라서 화면 위쪽이 뱃머리, 오른쪽이 우현
%   방향각 delta 는 x_b 에서 y_b 쪽으로 재므로 화면 극각 theta = 90 - delta
%
% 눈금
%   기하  12 mm/m  — b = 1 m 가 12 mm, L = 2 m 가 24 mm
%   힘    0.1 mm/N — 225 N 이 22.5 mm
%   (다)  0.3 mm/N — lambda 축
%
% 숫자의 출처 — A3_nullspace_example.m 실행 출력 (한 자리까지 같다)
%   요구 tau_c = [250; 100; 0],  손풀이 기하 b = 1 m, L = 2 m
%   보정 전  u_e  = [225;  50;  25;  50]   delta = 12.5288 / 63.4349 deg
%   보정 후  u_e* = [225;  75;  25;  25]   delta = 18.4349 / 45.0000 deg
%   lambda 구간  좌현 [-275, 175],  우현 [25, 75],  교집합 [25, 75],  lambda* = 25

\documentclass[border=6pt]{standalone}
% capstone-style.tex — 이 볼트의 TikZ 그림이 공유하는 스타일
%
% 쓰는 법:  % capstone-style.tex — 이 볼트의 TikZ 그림이 공유하는 스타일
%
% 쓰는 법:  % capstone-style.tex — 이 볼트의 TikZ 그림이 공유하는 스타일
%
% 쓰는 법:  \input{capstone-style}   (assets/src/ 안에서)
% 빌드:     bash _tools/tikz2svg.sh assets/src/w02-node-topic.tex assets/w02-node-topic.svg
%
% 왜 TikZ 인가
%   손으로 SVG 좌표를 쓰면 화살촉·선 굵기·글자 크기가 그림마다 미묘하게 달라진다.
%   그것이 "자료가 어설퍼 보인다"의 정체다. 스타일을 한 번만 정하고 상대좌표로
%   배치하면 좌표를 고쳐도 그림이 무너지지 않는다.
%   수식은 본문과 같은 TeX 글꼴로 나온다.
%
% 한글
%   ko.TeX(DVI 경로)를 쓴다. XeLaTeX 은 TikZ 그래픽을 PDF special 로 내보내는데
%   dvisvgm 이 그것을 받으려면 Ghostscript < 10.01 이 필요하다. 이 PC 에는
%   10.07 이 깔려 있어 그 경로가 막힌다. latex -> DVI -> dvisvgm 이 유일하게
%   동작하는 조합이다.

\usepackage[hangul]{kotex}
\usepackage{tikz}
\usepackage{amsmath}
\usepackage{xcolor}
\usetikzlibrary{arrows.meta, positioning, calc, fit, backgrounds, decorations.pathmorphing}

% ---- 색 --------------------------------------------------------------------
% 강조는 하나만. 나머지는 검정. 흑백 인쇄에서도 읽혀야 한다.
\definecolor{ink}{HTML}{1A1A1A}
\definecolor{soft}{HTML}{666666}
\definecolor{accent}{HTML}{7C3AED}
\definecolor{warn}{HTML}{B91C1C}

% 계층 색 — Simulink 모델의 블록 색과 같은 규칙을 쓴다
\definecolor{guid}{HTML}{1D4ED8}    % 유도
\definecolor{ctrl}{HTML}{EA580C}    % 제어
\definecolor{thr}{HTML}{D97706}     % 추진기
\definecolor{plant}{HTML}{16A34A}   % 운동모델
\definecolor{nav}{HTML}{7C3AED}     % 항법 · 신호처리
\definecolor{pathc}{HTML}{0D9488}   % 계획 경로 · 항적
\definecolor{errc}{HTML}{C2410C}    % 오차
\definecolor{flow}{HTML}{0369A1}    % 조류 · 바람

% ---- 선과 화살촉 -----------------------------------------------------------
\tikzset{
  sig/.style   = {ink, line width=0.5pt, -{Stealth[length=4pt, width=3pt]}},
  wire/.style  = {ink, line width=0.5pt},
  fb/.style    = {ink, line width=0.5pt, densely dashed,
                  -{Stealth[length=4pt, width=3pt]}},
  asig/.style  = {accent, line width=0.5pt, densely dashed,
                  -{Stealth[length=4pt, width=3pt]}},
  % 블록 — 흰 바탕, 직각, 검은 테두리. 색은 테두리에만 준다
  blk/.style   = {draw=ink, line width=0.55pt, fill=white,
                  minimum width=18mm, minimum height=8.5mm, inner sep=2pt, font=\small},
  gblk/.style  = {blk, draw=guid},
  cblk/.style  = {blk, draw=ctrl},
  tblk/.style  = {blk, draw=thr},
  pblk/.style  = {blk, draw=plant},
  nblk/.style  = {blk, draw=nav},
  % 합산점 · 분기점
  sum/.style   = {draw=ink, line width=0.55pt, fill=white, circle,
                  inner sep=0pt, minimum size=4.6mm, font=\scriptsize},
  dot/.style   = {ink, fill=ink, circle, inner sep=0pt, minimum size=1.5mm},
  % 글자
  sname/.style = {font=\small, inner sep=1.5pt},
  note/.style  = {font=\scriptsize, soft, inner sep=1.5pt},
  wnote/.style = {font=\scriptsize, warn, inner sep=1.5pt},
  anote/.style = {font=\scriptsize, accent, inner sep=1.5pt},
  axis/.style  = {ink, line width=0.5pt},
  tick/.style  = {font=\scriptsize, soft},
  % 묶음 상자 — 서브시스템 표시
  grp/.style   = {draw=soft, line width=0.4pt, densely dotted, rounded corners=1mm,
                  inner sep=3mm},
}
   (assets/src/ 안에서)
% 빌드:     bash _tools/tikz2svg.sh assets/src/w02-node-topic.tex assets/w02-node-topic.svg
%
% 왜 TikZ 인가
%   손으로 SVG 좌표를 쓰면 화살촉·선 굵기·글자 크기가 그림마다 미묘하게 달라진다.
%   그것이 "자료가 어설퍼 보인다"의 정체다. 스타일을 한 번만 정하고 상대좌표로
%   배치하면 좌표를 고쳐도 그림이 무너지지 않는다.
%   수식은 본문과 같은 TeX 글꼴로 나온다.
%
% 한글
%   ko.TeX(DVI 경로)를 쓴다. XeLaTeX 은 TikZ 그래픽을 PDF special 로 내보내는데
%   dvisvgm 이 그것을 받으려면 Ghostscript < 10.01 이 필요하다. 이 PC 에는
%   10.07 이 깔려 있어 그 경로가 막힌다. latex -> DVI -> dvisvgm 이 유일하게
%   동작하는 조합이다.

\usepackage[hangul]{kotex}
\usepackage{tikz}
\usepackage{amsmath}
\usepackage{xcolor}
\usetikzlibrary{arrows.meta, positioning, calc, fit, backgrounds, decorations.pathmorphing}

% ---- 색 --------------------------------------------------------------------
% 강조는 하나만. 나머지는 검정. 흑백 인쇄에서도 읽혀야 한다.
\definecolor{ink}{HTML}{1A1A1A}
\definecolor{soft}{HTML}{666666}
\definecolor{accent}{HTML}{7C3AED}
\definecolor{warn}{HTML}{B91C1C}

% 계층 색 — Simulink 모델의 블록 색과 같은 규칙을 쓴다
\definecolor{guid}{HTML}{1D4ED8}    % 유도
\definecolor{ctrl}{HTML}{EA580C}    % 제어
\definecolor{thr}{HTML}{D97706}     % 추진기
\definecolor{plant}{HTML}{16A34A}   % 운동모델
\definecolor{nav}{HTML}{7C3AED}     % 항법 · 신호처리
\definecolor{pathc}{HTML}{0D9488}   % 계획 경로 · 항적
\definecolor{errc}{HTML}{C2410C}    % 오차
\definecolor{flow}{HTML}{0369A1}    % 조류 · 바람

% ---- 선과 화살촉 -----------------------------------------------------------
\tikzset{
  sig/.style   = {ink, line width=0.5pt, -{Stealth[length=4pt, width=3pt]}},
  wire/.style  = {ink, line width=0.5pt},
  fb/.style    = {ink, line width=0.5pt, densely dashed,
                  -{Stealth[length=4pt, width=3pt]}},
  asig/.style  = {accent, line width=0.5pt, densely dashed,
                  -{Stealth[length=4pt, width=3pt]}},
  % 블록 — 흰 바탕, 직각, 검은 테두리. 색은 테두리에만 준다
  blk/.style   = {draw=ink, line width=0.55pt, fill=white,
                  minimum width=18mm, minimum height=8.5mm, inner sep=2pt, font=\small},
  gblk/.style  = {blk, draw=guid},
  cblk/.style  = {blk, draw=ctrl},
  tblk/.style  = {blk, draw=thr},
  pblk/.style  = {blk, draw=plant},
  nblk/.style  = {blk, draw=nav},
  % 합산점 · 분기점
  sum/.style   = {draw=ink, line width=0.55pt, fill=white, circle,
                  inner sep=0pt, minimum size=4.6mm, font=\scriptsize},
  dot/.style   = {ink, fill=ink, circle, inner sep=0pt, minimum size=1.5mm},
  % 글자
  sname/.style = {font=\small, inner sep=1.5pt},
  note/.style  = {font=\scriptsize, soft, inner sep=1.5pt},
  wnote/.style = {font=\scriptsize, warn, inner sep=1.5pt},
  anote/.style = {font=\scriptsize, accent, inner sep=1.5pt},
  axis/.style  = {ink, line width=0.5pt},
  tick/.style  = {font=\scriptsize, soft},
  % 묶음 상자 — 서브시스템 표시
  grp/.style   = {draw=soft, line width=0.4pt, densely dotted, rounded corners=1mm,
                  inner sep=3mm},
}
   (assets/src/ 안에서)
% 빌드:     bash _tools/tikz2svg.sh assets/src/w02-node-topic.tex assets/w02-node-topic.svg
%
% 왜 TikZ 인가
%   손으로 SVG 좌표를 쓰면 화살촉·선 굵기·글자 크기가 그림마다 미묘하게 달라진다.
%   그것이 "자료가 어설퍼 보인다"의 정체다. 스타일을 한 번만 정하고 상대좌표로
%   배치하면 좌표를 고쳐도 그림이 무너지지 않는다.
%   수식은 본문과 같은 TeX 글꼴로 나온다.
%
% 한글
%   ko.TeX(DVI 경로)를 쓴다. XeLaTeX 은 TikZ 그래픽을 PDF special 로 내보내는데
%   dvisvgm 이 그것을 받으려면 Ghostscript < 10.01 이 필요하다. 이 PC 에는
%   10.07 이 깔려 있어 그 경로가 막힌다. latex -> DVI -> dvisvgm 이 유일하게
%   동작하는 조합이다.

\usepackage[hangul]{kotex}
\usepackage{tikz}
\usepackage{amsmath}
\usepackage{xcolor}
\usetikzlibrary{arrows.meta, positioning, calc, fit, backgrounds, decorations.pathmorphing}

% ---- 색 --------------------------------------------------------------------
% 강조는 하나만. 나머지는 검정. 흑백 인쇄에서도 읽혀야 한다.
\definecolor{ink}{HTML}{1A1A1A}
\definecolor{soft}{HTML}{666666}
\definecolor{accent}{HTML}{7C3AED}
\definecolor{warn}{HTML}{B91C1C}

% 계층 색 — Simulink 모델의 블록 색과 같은 규칙을 쓴다
\definecolor{guid}{HTML}{1D4ED8}    % 유도
\definecolor{ctrl}{HTML}{EA580C}    % 제어
\definecolor{thr}{HTML}{D97706}     % 추진기
\definecolor{plant}{HTML}{16A34A}   % 운동모델
\definecolor{nav}{HTML}{7C3AED}     % 항법 · 신호처리
\definecolor{pathc}{HTML}{0D9488}   % 계획 경로 · 항적
\definecolor{errc}{HTML}{C2410C}    % 오차
\definecolor{flow}{HTML}{0369A1}    % 조류 · 바람

% ---- 선과 화살촉 -----------------------------------------------------------
\tikzset{
  sig/.style   = {ink, line width=0.5pt, -{Stealth[length=4pt, width=3pt]}},
  wire/.style  = {ink, line width=0.5pt},
  fb/.style    = {ink, line width=0.5pt, densely dashed,
                  -{Stealth[length=4pt, width=3pt]}},
  asig/.style  = {accent, line width=0.5pt, densely dashed,
                  -{Stealth[length=4pt, width=3pt]}},
  % 블록 — 흰 바탕, 직각, 검은 테두리. 색은 테두리에만 준다
  blk/.style   = {draw=ink, line width=0.55pt, fill=white,
                  minimum width=18mm, minimum height=8.5mm, inner sep=2pt, font=\small},
  gblk/.style  = {blk, draw=guid},
  cblk/.style  = {blk, draw=ctrl},
  tblk/.style  = {blk, draw=thr},
  pblk/.style  = {blk, draw=plant},
  nblk/.style  = {blk, draw=nav},
  % 합산점 · 분기점
  sum/.style   = {draw=ink, line width=0.55pt, fill=white, circle,
                  inner sep=0pt, minimum size=4.6mm, font=\scriptsize},
  dot/.style   = {ink, fill=ink, circle, inner sep=0pt, minimum size=1.5mm},
  % 글자
  sname/.style = {font=\small, inner sep=1.5pt},
  note/.style  = {font=\scriptsize, soft, inner sep=1.5pt},
  wnote/.style = {font=\scriptsize, warn, inner sep=1.5pt},
  anote/.style = {font=\scriptsize, accent, inner sep=1.5pt},
  axis/.style  = {ink, line width=0.5pt},
  tick/.style  = {font=\scriptsize, soft},
  % 묶음 상자 — 서브시스템 표시
  grp/.style   = {draw=soft, line width=0.4pt, densely dotted, rounded corners=1mm,
                  inner sep=3mm},
}


\newcommand{\drawhull}{%
  \draw[soft, line width=0.5pt]
       (-15.5,-28) -- (-15.5,24) -- (-12,31) -- (-8.5,24) -- (-8.5,-28) -- cycle;
  \draw[soft, line width=0.5pt]
       ( 15.5,-28) -- ( 15.5,24) -- ( 12,31) -- (  8.5,24) -- (  8.5,-28) -- cycle;
  \draw[soft, line width=0.5pt] (-8.5, 10) -- (8.5, 10);
  \draw[soft, line width=0.5pt] (-8.5,-10) -- (8.5,-10);
  \node[font=\scriptsize, soft, anchor=south] at (0,32) {선수};
}

\newcommand{\drawwedge}{%
  \fill[guid!10] (0,0) -- (45:15mm) arc (45:135:15mm) -- cycle;
  \draw[guid, line width=0.35pt] (45:15mm) arc (45:135:15mm);
  \draw[soft, line width=0.35pt, densely dashed] (0,0) -- (45:16.5mm);
  \draw[soft, line width=0.35pt, densely dashed] (0,0) -- (135:16.5mm);
}

\begin{document}
\begin{tikzpicture}[x=1mm, y=1mm]

\tikzset{
  fvec/.style = {thr,   line width=1.4pt, -{Stealth[length=4.2pt,width=3pt]}},
  bvec/.style = {warn,  line width=1.4pt, -{Stealth[length=4.2pt,width=3pt]}},
  ovec/.style = {soft,  line width=0.8pt, densely dashed,
                 -{Stealth[length=3.4pt,width=2.4pt]}},
  lam/.style  = {plant, line width=0.9pt, -{Stealth[length=3.6pt,width=2.6pt]}},
  axs/.style  = {nav,   line width=0.6pt, -{Stealth[length=3.4pt,width=2.4pt]}},
  hub/.style  = {circle, draw=thr, fill=white, line width=0.7pt,
                 minimum size=3mm, inner sep=0pt},
  tag/.style  = {circle, draw=ink, fill=white, line width=0.4pt,
                 minimum size=3.6mm, inner sep=0pt, font=\tiny},
}

% ===========================================================================
% (가) 보정 전
% ===========================================================================
\begin{scope}[shift={(0,0)}]

  \node[font=\small\bfseries, anchor=north west] at (0,56)
       {(가) 보정 전 — 우현 추진기가 쐐기를 넘는다};

  \begin{scope}[shift={(44,14)}]
    \drawhull

    % --- 좌현 추진기 ------------------------------------------------------
    \begin{scope}[shift={(-12,-24)}]
      \drawwedge
      \draw[fvec] (0,0) -- (77.4712:23.05mm);
      \node[tag] at (77.4712:25.5mm) {L};
      \node[hub] at (0,0) {};
    \end{scope}

    % --- 우현 추진기 — 쐐기 밖 -------------------------------------------
    \begin{scope}[shift={(12,-24)}]
      \drawwedge
      \draw[bvec] (0,0) -- (26.5651:5.59mm);
      \draw[warn, line width=0.4pt, densely dotted] (0,0) -- (26.5651:17mm);
      \node[tag, draw=warn] at (26.5651:19.5mm) {R};
      \node[font=\scriptsize, warn, anchor=south west, inner sep=1.5pt]
           at (26.5651:21mm) {쐐기 밖};
      \node[hub] at (0,0) {};
    \end{scope}

    % --- 숫자 ------------------------------------------------------------
    \node[anchor=north west, font=\scriptsize, align=left] at (-34,-34)
         {$\mathbf{u}_e = \mathbf{B}^{+}\boldsymbol{\tau}_c
           = [\,225,\ 50,\ 25,\ 50\,]^{\mathsf T}$ \\[3pt]
          L \ \ $T = 230.49$ N, \ $\delta = 12.53^\circ$ \quad
          \textcolor{plant}{쐐기 안} \\[2pt]
          R \ \ $T = 55.90$ N, \ $\delta = 63.43^\circ$ \quad
          \textcolor{warn}{쐐기 밖 — $18.43^\circ$ 초과}};
  \end{scope}

  \node[anchor=north west, font=\scriptsize, text width=88mm] at (0,-34)
       {\textbf{최소 노름 해는 각도 한계를 모른다} \\[4pt]
        $\mathbf{B}^{+}$ 는 $\lVert \mathbf{u}_e \rVert$ 만 가장 작게 만든다.
        쐐기는 계산에 들어가 있지 않다 \\[4pt]
        우현의 $F_x$ 가 $25$ N 밖에 안 되는데 $F_y$ 가 $50$ N 이라
        $\delta = \arctan(50/25) = 63.43^\circ$ 가 된다 \\[4pt]
        \textcolor{warn}{\textbf{여기서 그냥 깎으면 요구가 줄어든다}} \\[2pt]
        규칙 (c)가 $45^\circ$ 로 깎는다. 실제로 나는 렌치는
        $[\,264.53,\ 89.53,\ 6.41\,]$ 로, 요구와
        $\lVert \mathbf{e} \rVert = 19.02$ 만큼 어긋난다 \\[2pt]
        요구하지도 않은 선회 모멘트 $6.41$ N\,m 가 따라 나온다};
\end{scope}

% ===========================================================================
% (나) 보정 후
% ===========================================================================
\begin{scope}[shift={(98,0)}]

  \node[font=\small\bfseries, anchor=north west] at (0,56)
       {(나) 보정 후 — 두 힘을 $+\lambda$ 와 $-\lambda$ 로 반대로 민다};

  \begin{scope}[shift={(44,14)}]
    \drawhull

    % --- 좌현 — F_y 가 50 -> 75 (+lambda) --------------------------------
    \begin{scope}[shift={(-12,-24)}]
      \drawwedge
      \draw[ovec] (0,0) -- (77.4712:23.05mm);
      \draw[fvec] (0,0) -- (71.5651:23.72mm);
      \node[tag] at (71.5651:28mm) {L};
      % 미끄러짐 — 끝점이 오른쪽으로 lambda*0.1 = 2.5 mm
      \draw[lam] (5,22.5) -- (7.5,22.5);
      \node[font=\scriptsize, plant, anchor=north, inner sep=1.5pt] at (6.2,21.6)
           {$+\lambda$};
    \end{scope}

    % --- 우현 — F_y 가 50 -> 25 (-lambda) --------------------------------
    \begin{scope}[shift={(12,-24)}]
      \drawwedge
      \draw[ovec] (0,0) -- (26.5651:5.59mm);
      \draw[fvec] (0,0) -- (45:3.54mm);
      \draw[plant, line width=0.4pt, densely dotted] (0,0) -- (45:17mm);
      \node[tag, draw=plant] at (45:19.5mm) {R};
      \draw[lam] (5,2.5) -- (2.5,2.5);
      \node[font=\scriptsize, plant, anchor=north, inner sep=1.5pt] at (3.8,1.8)
           {$-\lambda$};
      \node[font=\scriptsize, plant, anchor=south west, inner sep=1.5pt]
           at (45:21mm) {한계에 붙음};
      \node[hub] at (0,0) {};
    \end{scope}
    \begin{scope}[shift={(-12,-24)}]
      \node[hub] at (0,0) {};
    \end{scope}

    % --- 숫자 ------------------------------------------------------------
    \node[anchor=north west, font=\scriptsize, align=left] at (-34,-34)
         {$\mathbf{u}_e^{\ast} = \mathbf{u}_e + 25\,\mathbf{n}
           = [\,225,\ 75,\ 25,\ 25\,]^{\mathsf T}$,
           \quad $\mathbf{n} = [\,0,\ 1,\ 0,\ -1\,]^{\mathsf T}$ \\[3pt]
          L \ \ $T = 237.17$ N, \ $\delta = 18.43^\circ$ \quad
          \textcolor{plant}{쐐기 안} \\[2pt]
          R \ \ $T = 35.36$ N, \ $\delta = 45.00^\circ$ \quad
          \textcolor{plant}{경계에 정확히 붙음}};
  \end{scope}

  \node[anchor=north west, font=\scriptsize, text width=88mm] at (0,-34)
       {\textbf{요구 렌치는 하나도 바뀌지 않는다} \ \
        $\mathbf{B}\mathbf{u}_e^{\ast} = [\,250,\ 100,\ 0\,]^{\mathsf T}$ \\[4pt]
        \hspace*{2mm}$X$ \ \ $F_x$ 를 건드리지 않았다 \\[2pt]
        \hspace*{2mm}$Y$ \ \ 한쪽은 $+\lambda$, 다른 쪽은 $-\lambda$ —
        더하면 $0$ \\[2pt]
        \hspace*{2mm}$N$ \ \ 두 추진기가 \textbf{같은 전후 위치}
        $x_b = -L$ 에 있다. 모멘트 팔이 둘 다 $-L$ 로 같으므로
        $-L(+\lambda) - L(-\lambda) = 0$ \\[4pt]
        \textcolor{plant}{\textbf{대가는 추력 노름뿐이다}} \ \
        $\lVert \mathbf{u}_e \rVert$ 가 $237.17 \to 239.79$ ($+1.11\,\%$).
        최소 노름 해에서 벗어났으니 당연하다 \\[2pt]
        $1.11\,\%$ 더 쓰고 못 낸 힘 $19.02$ 를 $0$ 으로 만든다};
\end{scope}

% ===========================================================================
% (다) lambda 를 고르는 법 — 1 차 부등식 둘
% ===========================================================================
\draw[wire, soft] (0,-84) -- (186,-84);

\node[font=\small\bfseries, anchor=north west] at (0,-89)
     {(다) $\lambda$ 는 부등식 둘의 교집합에서 고른다 — 반복도 최적화기도 없다};

\begin{scope}[shift={(95,-132)}]

  % --- 축 ----------------------------------------------------------------
  \draw[axs] (-95,0) -- (62,0);
  \node[font=\scriptsize, nav, anchor=west] at (63,0) {$\lambda$ [N]};

  \foreach \v/\x in {-275/-82.5, 0/0, 25/7.5, 75/22.5, 175/52.5} {
    \draw[ink, line width=0.4pt] (\x,-1.6) -- (\x,1.6);
  }
  \node[font=\scriptsize, soft, anchor=north] at (-82.5,-2) {$-275$};
  \node[font=\scriptsize, soft, anchor=north] at (52.5,-2)  {$175$};
  \node[font=\scriptsize, soft, anchor=north east, inner sep=1pt] at (-1,-2) {$0$};
  \node[font=\scriptsize, plant, anchor=north west, inner sep=1pt] at (8.5,-2) {$25$};
  \node[font=\scriptsize, plant, anchor=north west, inner sep=1pt] at (23,-2) {$75$};

  % --- 좌현 구간 ----------------------------------------------------------
  \draw[guid, line width=2pt] (-82.5,12) -- (52.5,12);
  \draw[guid, line width=0.8pt] (-82.5,9.5) -- (-82.5,14.5);
  \draw[guid, line width=0.8pt] ( 52.5,9.5) -- ( 52.5,14.5);
  \node[font=\scriptsize, guid, anchor=south west, inner sep=1.5pt] at (-82.5,15)
       {좌현 \ $\lvert 50 + \lambda \rvert \le \tan 45^\circ \cdot \lvert 225 \rvert$
        \quad $\to$ \quad $-275 \le \lambda \le 175$ \ \ (넉넉하다)};

  % --- 우현 구간 ----------------------------------------------------------
  \draw[ctrl, line width=2pt] (7.5,27) -- (22.5,27);
  \draw[ctrl, line width=0.8pt] ( 7.5,24.5) -- ( 7.5,29.5);
  \draw[ctrl, line width=0.8pt] (22.5,24.5) -- (22.5,29.5);
  \node[font=\scriptsize, ctrl, anchor=south west, inner sep=1.5pt] at (-82.5,30)
       {우현 \ $\lvert 50 - \lambda \rvert \le \tan 45^\circ \cdot \lvert 25 \rvert$
        \quad $\to$ \quad $25 \le \lambda \le 75$ \ \
        \textbf{여기가 답을 결정한다}};
  \draw[ctrl, line width=0.35pt] (-20,30) -- (7,28);

  % --- 교집합 -------------------------------------------------------------
  \draw[plant, line width=2.6pt] (7.5,-14) -- (22.5,-14);
  \node[font=\scriptsize, plant, anchor=west, inner sep=1.5pt] at (24,-14)
       {교집합 $[\,25,\ 75\,]$ — 이 안이면 두 추진기 모두 쐐기 안};

  % --- lambda = 0 ---------------------------------------------------------
  \draw[errc, line width=0.6pt, densely dashed] (0,-22) -- (0,20);
  \node[font=\scriptsize, errc, anchor=north east, inner sep=1.5pt, align=right]
       at (-1,-23)
       {$\lambda = 0$ — 보정하지 않음 \\ 교집합 밖이다};

  % --- lambda* = 25 -------------------------------------------------------
  \draw[plant, line width=0.7pt] (7.5,-18) -- (7.5,24);
  \node[circle, fill=plant, inner sep=0pt, minimum size=1.8mm] at (7.5,-14) {};
  \node[font=\scriptsize, plant, anchor=north west, inner sep=1.5pt, align=left]
       at (9,-23)
       {$\lambda^{\ast} = 25$ — 교집합 안에서 $0$ 에 가장 가까운 값 \\
        \texttt{lam = min(max(0, lo), hi)} 한 줄이 이것을 고른다};
\end{scope}

\node[anchor=north west, font=\scriptsize, text width=88mm] at (0,-172)
     {\textbf{부등식이 1 차인 이유} \\[4pt]
      각도 한계 $\lvert \delta_i \rvert \le \delta_{\max}$ 를 그대로 쓰면
      $\arctan$ 이 들어가 풀기 어렵다 \\[2pt]
      양변에 $\tan$ 을 씌우면
      $\lvert F_{y,i} \rvert \le \tan\delta_{\max}\,\lvert F_{x,i} \rvert$ —
      $\lambda$ 에 대한 1 차 부등식이 된다 \\[2pt]
      $F_x$ 는 $\mathbf{n}$ 이 건드리지 않으므로 우변이 상수다};

\node[anchor=north west, font=\scriptsize, text width=88mm] at (98,-172)
     {\textbf{교집합이 비면} \\[4pt]
      어떤 $\lambda$ 로도 두 추진기를 동시에 쐐기 안에 넣을 수 없다는 뜻이다 \\[2pt]
      이때는 가운데 값 $(\mathrm{lo}+\mathrm{hi})/2$ 를 쓴다 —
      어느 한쪽만 맞추는 것보다 \textbf{덜 넘긴다} \\[4pt]
      \textcolor{warn}{\textbf{보정은 못 내는 요구를 내게 해 주지 않는다}} \\[2pt]
      쐐기 밖의 요구를 쐐기 안으로 \textbf{옮겨 줄 뿐}이다.
      요구 자체가 추력 한계를 넘으면 규칙 (d)가 그대로 깎는다};

\end{tikzpicture}
\end{document}
